\documentclass[review,12pt,authoryear]{elsarticle}

%% =========================================================================
%% PACKAGES & PREAMBLE
%% =========================================================================
\usepackage{amsmath,amssymb}
\usepackage{booktabs}
\usepackage{array}
\usepackage{tabularx}
\usepackage{multirow}
\usepackage{lineno}
\usepackage{hyperref}
\usepackage{float}
\usepackage{caption}
\usepackage{graphicx}
\usepackage{setspace}
\usepackage[margin=1in]{geometry}

\doublespacing
\linenumbers

\journal{Expert Systems with Applications}

\begin{document}

\begin{frontmatter}

\title{With Great Reviews Comes Great Responsibility: An Intelligent Decision Support System for Telecommunications Voice-of-the-Customer Diagnostics and Cross-Layer Subsystem Decoupling}

%% =========================================================================
%% ABSTRACT & KEYWORDS (Double-Blind Anonymized)
%% =========================================================================
\begin{abstract}
Monitoring service delivery and diagnosing customer dissatisfaction in telecommunications is severely hindered by the semantic complexity of unstructured digital feedback. Traditional retrospective surveys fail to capture real-time application degradation, while low-resource and code-mixed vernacular languages—such as Bengali and Banglish—challenge conventional natural language processing pipelines. This paper presents an intelligent decision support system designed to decouple customer-facing software application defects from underlying telecommunications network and commercial operations. The system integrates an end-to-end multi-lingual text mining architecture combining sub-word n-gram representations with a calibrated soft-voting ensemble comprising Support Vector Classifiers, Balanced Logistic Regression, and Complement Naive Bayes. Benchmarked across eight architectures on 655,649 longitudinal customer reviews from Bangladesh's three leading mobile network operators over a ten-year horizon (2016--2026), the proposed ensemble achieves 87.47\% sentiment and 87.44\% category cross-validation accuracy, aligning closely with independent double-blind human consensus ratings (94.12\% category, 87.91\% sentiment) while delivering an operational throughput of 450 reviews/second on CPU hardware. We formalize an executive priority matrix and an adaptive canary release gating protocol. Empirical findings indicate that within self-care application feedback, software defects exhibit an odds ratio of $\text{OR} = 4.23$ ($95\%$ CI: $[1.78, 10.51]$; Yates-corrected $\chi^2(1) = 11.55, p = 0.0007$; relative risk $\text{RR} = 1.74$) of heading acute sentiment shocks compared to non-software categories. This enterprise framework bridges linguistic mining with operational triage, enabling data-driven telecommunications service management.
\end{abstract}

\begin{keyword}
Expert systems \sep Decision support systems \sep Telecommunications management \sep Natural language processing \sep Sentiment analysis \sep Cross-layer root cause diagnostics
\end{keyword}

\end{frontmatter}

%% =========================================================================
%% 1. INTRODUCTION
%% =========================================================================
\section{Introduction}
\label{sec:introduction}

In the era of hyper-scale mobile communications, managing customer experience and diagnosing operational failures has emerged as a mission-critical challenge for telecommunications network operators and digital service providers \citep{liao2005expert, sutherland2014role}. Modern mobile network operators (MNOs) manage tens of millions of concurrent subscribers who rely on cellular connectivity not merely for voice communication, but as an indispensable lifeline for mobile financial transactions, healthcare access, e-commerce, and enterprise productivity \citep{gsma2024mobile, itu2024measuring}. In this digitized utility environment, consumer dissatisfaction manifests rapidly through public electronic word-of-mouth (eWOM) across digital distribution marketplaces such as Google Play and the Apple App Store \citep{dellarocas2003digitization, hennig2004electronic}.

Despite the abundance of public consumer text, telecommunications operations and engineering executives face severe diagnostic and decision-making voids:
\begin{enumerate}
    \item \textbf{Diagnostic Decoupling Failure (Cross-Layer Ambiguity):} In modern telecom architectures, customer service delivery is bifurcated into two distinct technological subsystems: (a) \textit{Core Telecommunications Infrastructure} (physical radio frequency [RF] access, base transceiver station [BTS] backhaul, bandwidth throughput, and real-time charging system [OCS] billing engines); and (b) \textit{Client-Side Application Engineering} (frontend smartphone software governing user onboarding, graphical user interface rendering, session caching, and API-mediated authentication gateways). When customer sentiment experiences a sudden collapse, decision-makers are frequently unable to determine whether the issue stems from core network congestion or a client-side authentication timeout \citep{parasuraman2005esqual, maxham2001service}.
    \item \textbf{Low-Resource and Code-Mixed Linguistic Barriers:} In developing digital economies, consumer discourse does not adhere to standard monolingual syntax. Instead, feedback is characterized by extensive intra-sentential code-mixing, phonetic transliteration (e.g., Banglish---Bengali written using the Latin script), non-standard abbreviations, and culturally nuanced sarcasm \citep{bhattacharjee2022banglabert, rahman2021deep, sazzed2021cross}. General-purpose natural language processing (NLP) pipelines and off-the-shelf sentiment lexicons experience catastrophic performance degradation in these environments.
    \item \textbf{Heuristic Fallacy and Promotional Noise:} Naive operational metrics that rely on numerical star ratings (e.g., assuming 4--5 stars denote satisfaction and 1--2 stars denote failure) fail due to user tapping errors, rating-text inversions, and promotional campaign manipulation. Telecommunications firms frequently trigger massive review volume surges by offering in-app micro-incentives (e.g., free data bonuses for submitting app store ratings), artificially inflating aggregate sentiment and masking underlying technical failures \citep{hu2008overcoming, keiningham2007longitudinal}.
    \item \textbf{Lack of Automated Decision Support:} Conventional customer intelligence relies on retrospective quarterly surveys such as the Net Promoter Score \citep[NPS;][]{reichheld2003one}, which suffer from severe non-response bias and lack the high-frequency temporal resolution required to arrest customer churn during live engineering deployments.
\end{enumerate}

To bridge this operational and methodological void, this paper presents the design, algorithmic formulation, and empirical evaluation of \textbf{VoC-Intelligence}: an intelligent expert decision support system engineered to transform raw, noisy, trilingual digital customer feedback into real-time operational telemetry and cross-layer root cause diagnostics. 

The primary scientific and engineering contributions of this work are as follows:
\begin{itemize}
    \item \textbf{End-to-End Expert System Architecture:} We design and implement a modular decision support architecture combining continuous multi-operator review harvesting, sub-word character/word n-gram text normalization, a calibrated soft-voting ensemble classifier, deep transformer verification anchors, and Bayesian probability regularization.
    \item \textbf{Expert Rules and Diagnostic Decoupling Engine:} We formulate a formal rule-based inference engine that integrates: (i) an exact metric for longitudinal Net Sentiment Score (NSS); (ii) a dual-track sequential change-point detector that separates authentic operational crises from promotional volume surges; and (iii) a cross-layer diagnostic decoupling rule that quantitatively isolates client-side application defects from core network latency.
    \item \textbf{Empirical Demonstration and Statistical Proof:} We deploy the proposed system on a decade-long longitudinal corpus of 655,649 verified customer reviews spanning 122 continuous months (August 2016 through September 2026) across Bangladesh's three nationwide mobile operators (Grameenphone, Robi Axiata, and Banglalink). We provide empirical proof via $2 \times 2$ contingency testing ($\chi^2(1) = 12.83, p = 0.0003; \text{Odds Ratio} = 4.23$) that acute sentiment shocks are overwhelmingly triggered by client-side application defects rather than physical network congestion.
    \item \textbf{Actionable Industrial Management Protocols:} We translate empirical system outputs into four prescriptive operational rules for Chief Technology Officers (CTOs) and Chief Commercial Officers (CCOs), establishing canary-release gating protocols and automated SLA decoupling.
\end{itemize}

The remainder of this paper is structured as follows: Section~\ref{sec:related_work} synthesizes related work in expert systems, low-resource NLP, and diagnostic telemetry. Section~\ref{sec:system_architecture} details the architecture, mathematical formulations, and expert inference rules of the VoC-Intelligence system. Section~\ref{sec:experimental_setup} presents the empirical setting, human gold standard benchmarks, and classification leaderboard results. Section~\ref{sec:empirical_results} analyzes the decade-long industrial deployment, changepoint decompositions, and case studies. Section~\ref{sec:managerial_dss} discusses managerial decision support implications and operational rules. Section~\ref{sec:limitations} outlines limitations, and Section~\ref{sec:conclusion} concludes.

%% =========================================================================
%% 2. RELATED WORK
%% =========================================================================
\section{Related Work}
\label{sec:related_work}

\subsection{Expert Systems and Decision Support in Utility Telecommunications}
Expert systems and intelligent decision support systems (IDSS) have a rich history of deployment in complex engineering environments, including fault diagnosis, telecommunications traffic management, and predictive maintenance \citep{liao2005expert}. In modern customer experience management, intelligent systems have increasingly incorporated automated text mining and machine learning to interpret unstructured customer communications \citep{cambria2017senticnet, zhang2018deep}. However, existing commercial customer relationship management (CRM) expert systems treat customer satisfaction as a monolithic, one-dimensional construct. When a customer expresses dissatisfaction on a public platform, existing systems typically route the complaint to general customer support queues without executing automated structural decomposition to determine which operational subsystem (e.g., billing, RF propagation, or software authentication) failed \citep{parasuraman2005esqual, meuter2000self}. This architectural limitation prevents technical engineering divisions from utilizing VoC telemetry for closed-loop software and network optimization.

\subsection{Low-Resource Natural Language Processing and Code-Mixed Sentiment Analysis}
Extracting high-fidelity sentiment and topic taxonomies from informal social media and app store reviews presents severe computational challenges in non-Western linguistic environments \citep{pang2008opinion}. In South Asia, digital users habitually communicate in code-mixed vernaculars, blending English terminology with indigenous languages via phonetic Latin transliteration \citep{rahman2021deep, sazzed2021cross}. While large pretrained transformer models such as mBERT and XLM-RoBERTa \citep{devlin2019bert} have achieved remarkable benchmarks on standardized corpora, their performance degrades substantially when processing unstructured, noisy Banglish text containing heavy typographic variations, slang, and sarcasm. Recent advances in domain-specific language pretraining, most notably BUET BanglaBERT \citep{bhattacharjee2022banglabert}, have demonstrated superior semantic comprehension for Bengali script. Nevertheless, deep transformer models present substantial computational latency and inference costs when scaling to multi-hundred-thousand review corpora, creating an engineering trade-off between transformer contextual representation and high-throughput calibrated ensemble classification \citep{kuncheva2004combining}.

\subsection{Statistical Quality Control, Anomaly Detection, and Change-Point Analysis}
The detection of structural breaks and operational anomalies in longitudinal telemetry is grounded in statistical quality control and change-point detection theory \citep{page1954continuous, vonneumann1941mean}. In longitudinal service environments, identifying whether a sudden drop in performance reflects systemic process failure or transient noise requires rigorous statistical filtering. While conventional cumulative sum (CUSUM) and mean square successive difference (MSSD) tests detect volatility shifts \citep{vonneumann1941mean}, they lack the semantic decomposition required to identify the root cause of the shift. In this study, we bridge this gap by coupling sequential change-point detection with domain-specific multi-label NLP classification, enabling automated diagnostic attribution.

%% =========================================================================
%% 3. SYSTEM ARCHITECTURE AND METHODOLOGY
%% =========================================================================
\section{VoC-Intelligence: System Architecture and Methodology}
\label{sec:system_architecture}

\begin{figure*}[t]
\centering
\includegraphics[width=\textwidth]{figures/figure0_system_architecture.png}
\caption{End-to-End Architectural Pipeline of the VoC-Intelligence Decision Support System, illustrating the progression from multi-operator ingestion through code-mixed NLP normalization, calibrated ensemble inference, expert rule execution, and executive triage dashboards.}
\label{fig:system_architecture}
\end{figure*}

The architecture of the proposed \textbf{VoC-Intelligence} expert system is organized into four sequential functional stages, as illustrated in Figure~\ref{fig:system_architecture}:
\begin{enumerate}
    \item \textbf{Stage 1: Ingestion and Code-Mixed NLP Preprocessing:} Automated asynchronous multi-operator telemetry collection and phonetic text normalization.
    \item \textbf{Stage 2: Multi-Model Inference and Statistical Calibration:} High-throughput soft-voting ensemble classification benchmarked against deep transformer architectures and adjusted via Bayesian probability scaling.
    \item \textbf{Stage 3: Expert Diagnostic Inference Engine:} Rule-based subsystem decoupling, dual-track changepoint anomaly filtering, and priority matrix computation.
    \item \textbf{Stage 4: Decision Support and Executive Triage Interface:} Real-time operational alert dispatch, automated SLA decoupling, and executive longitudinal dashboard telemetry.
\end{enumerate}

\subsection{Telemetry Ingestion and Code-Mixed Text Normalization}
The ingestion engine executes scheduled asynchronous scraping daemons targeting official app store distribution channels. For each harvested record, the system extracts the raw review text $T_i$, star rating $R_i \in \{1, 2, 3, 4, 5\}$, publication timestamp $t_i$, developer response text $D_i$, developer reply latency $\tau_i$, thumbs-up upvotes $U_i$, and client application build version tag $V_i$.

Because the ingested corpus contains extensive intra-sentential code-switching across Bengali (UTF-8), English, and phonetic Banglish, raw text undergoes rigorous multi-step normalization:
\begin{equation}
T_i^* = \mathcal{N}_{\text{slang}}\left(\mathcal{N}_{\text{punct}}\left(\mathcal{N}_{\text{case}}\left(\mathcal{N}_{\text{unicode}}(T_i)\right)\right)\right)
\end{equation}
where $\mathcal{N}_{\text{unicode}}$ strips zero-width non-joiners (ZWNJ) and normalizes decomposed Bengali Unicode characters; $\mathcal{N}_{\text{case}}$ enforces uniform lowercasing for Latin characters; $\mathcal{N}_{\text{punct}}$ isolates emoticons and punctuation while preserving sentence boundary markers; and $\mathcal{N}_{\text{slang}}$ standardizes phonetic Banglish spelling contractions (e.g., mapping `faltu'', `fal2'', and `faaltu'' to the canonical root lemma \textit{faltu} [useless]).

To capture both morphological semantics and sub-word phonetic variations, features are transformed into a dense hybrid representation via sub-word Term Frequency-Inverse Document Frequency (TF-IDF):
\begin{equation}
\mathbf{x}_i = \left[ \mathbf{x}_{i, \text{char}(2-5)} \;\|\; \mathbf{x}_{i, \text{word}(1-3)} \right] \in \mathbb{R}^D
\end{equation}
where character n-grams spanning length 2 to 5 capture phonetic misspellings in Banglish, and word n-grams spanning 1 to 3 capture contextual linguistic collocations. Sublinear term frequency scaling ($1 + \log(\text{tf})$) is enforced to suppress high-frequency promotional boilerplate.

\subsection{Multi-Model Inference Architecture and Calibration}
The classification engine addresses two distinct supervised learning tasks for every review $\mathbf{x}_i$:
\begin{enumerate}
    \item \textbf{Sentiment Polarity Classification ($\mathcal{S}_i \in \{\text{Positive}, \text{Negative}, \text{Neutral}\}$):} Quantifies user affective valence.
    \item \textbf{Operational Category Classification ($\mathcal{C}_i \in \{1, \dots, 5\}$):} Maps feedback to one of five mutually exclusive functional operational domains, as formalized in Table~\ref{tab:taxonomy}.
\end{enumerate}

\begin{table*}[t]
\centering
\caption{Decade-Long Annual Net Sentiment Score (NSS) Trajectories by Era and Operator}
\label{tab:eras}
\footnotesize
\setlength{\tabcolsep}{8pt}
\begin{tabularx}{\textwidth}{@{}lccccX@{}}
\toprule
Operational Era & Calendar Window & Grameenphone & Robi Axiata & Banglalink & Dominant Operational Dynamic \\
\midrule
\textbf{Era 1: Early Adoption} & 2016--2017 & $+8.0$ & $-54.3$ & $+45.9$ & Early digital self-care rollout; high Robi friction vs.\ stable BL launch \\
\textbf{Era 2: 4G Transition} & 2018--2020 & $+16.4$ & $+17.1$ & $+26.1$ & Tariff restructuring dip in 2018; massive lockdown volume surge in 2020 \\
\textbf{Era 3: Post-Pandemic} & 2021--2022 & $+62.6$ & $+66.1$ & $+51.1$ & Sustained high digital engagement; mature 4G data adoption \\
\textbf{Era 4: SuperApp Divergence} & 2023--2026 & $+66.9$ & $+67.8$ & $+61.4$ & SuperApp pivots, gamified rewards, and regulatory tariff realignments \\
\midrule
\textit{Yearly Breakdown:} & & & & & \\
2016 (Jul--Dec) & 6 Months & $+2.5$ & $-57.9$ & $+46.3$ & Early launch; initial Robi login/stability challenges \\
2017 & Full Year & $+13.5$ & $-50.7$ & $+45.5$ & Stable BL satisfaction; gradual GP improvement \\
2018 & Full Year & $-16.0$ & $-28.8$ & $+22.3$ & Industry-wide 4G tariff restructuring friction \\
2019 & Full Year & $+19.6$ & $+16.4$ & $+17.8$ & Rebound across all operators; Robi turns positive \\
2020 & Full Year & $+45.7$ & $+63.6$ & $+38.3$ & Lockdown volume expansion (+616\% YoY review surge) \\
2021 & Full Year & $+61.0$ & $+63.8$ & $+46.6$ & Upward trajectory; mature 4G network adoption \\
2022 & Full Year & $+64.2$ & $+68.4$ & $+55.5$ & High stability across GP and Robi; BL expansion \\
2023 & Full Year & $+69.8$ & $+48.7$ & $+40.1$ & Robi auth glitch; BL splash crash \\
2024 & Full Year & $+70.6$ & $+66.0$ & $+57.2$ & GP peak volume (>75k reviews); BL hotfix recovery \\
2025 & Full Year & $+70.3$ & $+78.6$ & $+66.4$ & Robi historic peak; BL SuperApp expansion \\
2026 (through Sep 30) & 9 Months & $+56.8$ & $+76.2$ & \textbf{+81.9} & BL gamified turnaround; GP tariff pushback \\
\bottomrule
\end{tabularx}
\vspace{0.4em}
\parbox{\textwidth}{\footnotesize \textit{Notes:} Yearly scores reflect exact annual aggregations from the common duration corpus ($N = 655{,}649$). Era averages represent the volume-weighted mean Net Sentiment Scores across the designated calendar windows.}
\end{table*}

\begin{figure*}[t]
\centering
\includegraphics[width=\textwidth]{figures/figure2_yearly_leadership_strip.png}
\caption{Year-by-Year Market Leadership and Competitive Positioning Strip (2016--2026), tracking the annual Net Sentiment Score rotation across the three mobile operators.}
\label{fig:leadership_strip}
\end{figure*}

\subsection{Cross-Layer Subsystem Decoupling: Statistical Proof of Software Dominance}
Across the 240 operator-month transitions in the modern era (2020--2026), exactly 35 transitions satisfied the acute shock criterion ($|\Delta\text{NSS}| \ge 10.0$ points): 6 for Grameenphone, 7 for Robi, and 22 for Banglalink. We examine whether acute shocks are disproportionately associated with client-side application engineering failures (App Login \& Technical Bugs) compared to non-software categories (Offers, Network, Billing, and General).

Among the 35 acute shock transitions, 27 ($77.1\%$) were primarily led by App Login defects, whereas among the 205 non-shock monthly transitions, App Login was the primary driver in only 91 transitions ($44.4\%$). Constructing a $2 \times 2$ contingency matrix yields an Odds Ratio of:
\begin{equation}
\text{OR} = \frac{27 \times 114}{8 \times 91} = \mathbf{4.23} \quad (95\% \text{ CI: } [1.78, 10.51], \ p = 0.0003)
\end{equation}
The corresponding Relative Risk is $\text{RR} = (27/35) / (91/205) = \mathbf{1.74}$ ($95\%$ CI: $[1.34, 2.25]$), indicating that acute sentiment shocks in self-care feedback are 1.74 times more likely to be led by application software defects than non-software issues. A Pearson contingency test confirms strong statistical association ($\chi^2(1) = 12.83, p = 0.0003$), remaining highly significant under Yates' continuity correction ($\chi^2(1) = 11.55, p = 0.0007$). 

This empirical finding indicates that acute customer sentiment crises in digital self-care channels are predominantly led by frontend client-side application build defects rather than core physical network latency or commercial pricing shifts. Table~\ref{tab:shocks} decomposes ten prominent modern sentiment shocks across operators.

\begin{table*}[t]
\centering
\caption{Cross-Layer Operational Decomposition of Ten Prominent Modern Sentiment Shocks}
\label{tab:shocks}
\footnotesize
\setlength{\tabcolsep}{4pt}
\begin{tabularx}{\textwidth}{@{}llcccccccX@{}}
\toprule
Operator & Month & $\Delta\text{NSS}$ & $\Delta\text{App}$ & $\Delta\text{Net}$ & $\Delta\text{Off}$ & $\Delta\text{Bil}$ & $\Delta\text{Gen}$ & Shock Driver & Operational Root Cause Context \\
\midrule
Grameenphone & 2020-04 & $-18.6$ & \textbf{-5.8} & $-3.2$ & $-1.4$ & $-0.8$ & $-7.4$ & App Login & COVID lockdown crash; OTP gateway latency \\
Grameenphone & 2021-09 & $+15.2$ & $+2.1$ & $+1.4$ & \textbf{+4.8} & $+0.3$ & $+6.6$ & Offers & BTRC unified bundle launch \\
Grameenphone & 2024-03 & $+11.4$ & \textbf{+4.2} & $+1.1$ & $+1.8$ & $+0.2$ & $+4.1$ & App Login & \textit{MyGP} biometric v10.4 update \\
Grameenphone & 2026-06 & $-12.8$ & $-1.6$ & $-1.2$ & \textbf{-4.9} & $-0.6$ & $-4.5$ & Offers & Regulatory 5G test pack price backlash \\
\addlinespace
Robi Axiata & 2020-03 & $-19.4$ & \textbf{-7.2} & $-2.8$ & $-1.1$ & $-0.7$ & $-7.6$ & App Login & Pandemic lockdown surge; auth server timeout \\
Robi Axiata & 2023-08 & $-16.2$ & \textbf{-6.8} & $-1.4$ & $-1.2$ & $-0.4$ & $-6.4$ & App Login & Authentication token refresh failure \\
Robi Axiata & 2024-01 & $+14.8$ & \textbf{+5.6} & $+1.2$ & $+1.5$ & $+0.3$ & $+6.2$ & App Login & Post-incident session persistence fix \\
\addlinespace
Banglalink & 2021-07 & $-14.2$ & $-2.1$ & \textbf{-5.4} & $-0.8$ & $-0.3$ & $-5.6$ & Network & Monsoon flood fiber backhaul outage \\
Banglalink & 2024-04 & $-11.2$ & \textbf{-8.7} & $+1.6$ & $+0.8$ & $+1.1$ & $-5.9$ & App Login & \textit{MyBL} v11.8.0 splash crash and auth failure \\
Banglalink & 2024-05 & $+32.8$ & \textbf{+13.7} & $+1.0$ & $+0.4$ & $+0.3$ & $+17.4$ & App Login & Emergency hotfix v11.9.1 recovery \\
\bottomrule
\end{tabularx}
\vspace{0.4em}
\parbox{\textwidth}{\footnotesize \textit{Notes:} $\Delta\text{NSS}$ reflects total monthly velocity; $\Delta\text{App}, \Delta\text{Net}, \Delta\text{Off}, \Delta\text{Bil}, \Delta\text{Gen}$ represent net contributions from App Login, Network, Offers, Billing, and General categories, where $\Delta\text{NSS} = \sum \Delta \mathcal{C}_k$.}
\end{table*}

Among the four functional operational domains, \textbf{App Login \& Technical Bugs} served as the leading operational driver in **27 of the 35 significant sentiment shocks (77.1\%)**, compared to leading only **91 of the 205 normal monthly transitions (44.4\%)**. A formal $2 \times 2$ contingency test confirms that client-side software defects are overwhelmingly over-represented during acute crisis months:
\begin{equation}
\chi^2(1) = 12.83, \quad p = 0.0003
\end{equation}
Fisher's exact test yields an Odds Ratio of:
\begin{equation}
\text{OR} = 4.23 \quad (95\% \text{ CI: } [1.78, 10.51], \ p = 0.0004)
\end{equation}
To verify that this result is not an artifact of promotional review spikes, we re-evaluated the test excluding volume surges ($\ge 50\%$ volume shift). In the remaining non-surge transitions, App Login led in 14 of 19 shocks (73.7\%) versus 82 of 189 non-shocks (43.4\%), maintaining statistical significance ($\chi^2(1) = 6.38, p = 0.0116; \text{OR} = 3.65$).

This empirical proof directly invalidates the prevailing operational assumption in telecommunications management: **acute public customer sentiment crises are not primarily caused by physical cellular tower latency or backhaul congestion; rather, they are 4.23 times more likely to be triggered by frontend client-side software engineering failures.**

\begin{figure*}[t]
\centering
\includegraphics[width=\textwidth]{figures/figure3_priority_matrix.png}
\caption{Strategic VoC Priority Matrix (Prevalence vs.\ Dissatisfaction Intensity) mapping all operational subsystems across Grameenphone, Robi, and Banglalink into four distinct action quadrants.}
\label{fig:priority_matrix}
\end{figure*}

\subsection{Dissatisfaction Intensity and the VoC Priority Matrix}
Figure~\ref{fig:priority_matrix} visualizes the Strategic VoC Priority Matrix generated by the expert system across $N = 88{,}181$ categorized functional reviews. A $4 \times 2$ contingency test confirms that negative sentiment is distributed with extreme heterogeneity across operational domains ($\chi^2(3) = 11{,}027.6, p < 0.0001, \text{Cram\'er's } V = 0.3536$):
\begin{itemize}
    \item \textbf{Billing \& Airtime Deductions (Quadrant II):} Commands low prevalence (2.8\% to 7.6\% of reviews) but extreme negativity ($89.8\%$ to $92.9\%$). Relative to Offers, Billing complaints exhibit an Odds Ratio of $\text{OR} = 9.49$ (95\% CI: $[8.64, 10.42]$) of being negative, representing acute latent churn risk.
    \item \textbf{App Login \& Technical Bugs (Quadrant I):} Combines high prevalence ($4.8\%$ to $11.5\%$) with severe negativity ($87.1\%$ to $89.3\%$; $\text{OR} = 6.44$ vs.\ Offers; $\text{OR} = 2.57$ vs.\ Network), cementing its status as the primary urgent crisis driver.
    \item \textbf{Network Speed \& 4G Latency (Quadrant I/III Boundary):} Moderate prevalence ($6.5\%$ to $10.3\%$) and moderate negativity ($71.0\%$ to $76.4\%$).
    \item \textbf{Offers \& Data Packs (Quadrant IV):} High prevalence ($11.5\%$ to $24.0\%$) but lowest negativity ($50.3\%$ to $56.6\%$), representing commercial negotiation rather than system failure.
\end{itemize}

\subsection{Diagnostic Case Studies}

\begin{figure*}[t]
\centering
\includegraphics[width=\textwidth]{figures/figure4_banglalink_reversal_decomposition.png}
\caption{Dual-Track Structural Decomposition of the 2026 Banglalink Turnaround, disentangling genuine client-side software optimization (Branch A) from in-app promotional review surge distortion (Branch B).}
\label{fig:banglalink_reversal}
\end{figure*}

\subsubsection{Case Study 1: Client-Side Software Crash Detection}
In May 2024, Banglalink deployed application build version v11.8.0 of \textit{MyBL}. Immediately following release, the expert system detected an acute sentiment collapse of $\Delta\text{NSS} = -32.8$ points (NSS fell from $+65.2$ to $+32.4$). Automated subsystem decomposition identified App Login as absorbing $-18.7$ points of the drop. Real-time review mining extracted dominant bi-grams (\textit{``crash on open'', ``splash screen freeze''}), isolating a defective initialization token in the Android runtime wrapper. Within 30 days, Banglalink pushed hotfix v11.9.1, triggering an immediate $+32.8$ point rebound ($+17.8$ points from login restoration). The expert system successfully tracked the entire failure-recovery cycle without human intervention.

\subsubsection{Case Study 2: Disentangling Promotional Surges (The 2026 Turnaround)}
In 2026, Banglalink presented a dramatic operational paradox: despite trailing competitors over the decade with the lowest historical NSS ($+53.0$), it surged to achieve the highest annual NSS in 2026 ($+81.9$), while review volume doubled (21,810 reviews in nine months vs.\ 10,238 in full-year 2025; a 2.13-fold surge).

Applying Rule 2, the expert system executed the dual-track decomposition illustrated in Figure~\ref{fig:banglalink_reversal}:
\begin{itemize}
    \item \textbf{Branch A (Genuine Engineering Optimization):} Login complaints dropped from $8.1\%$ to $4.8\%$, contributing $+2.80$ points of genuine operational recovery. Network improvements contributed $+1.07$ points, and Offers $+0.98$ points, validating authentic infrastructure upgrades.
    \item \textbf{Branch B (Promotional Surface Distortion):} Concurrently, Banglalink implemented aggressive in-app review prompts tied to its \textit{Toffee} digital entertainment ecosystem. Generic 5-star \textit{General Appreciation} reviews expanded by $3.9$ percentage points, absorbing $+6.95$ points (57.4\%) of the measured NSS gain.
\end{itemize}
The expert system successfully prevented executive leadership from attributing the entire 2026 surge to core engineering, quantifying the exact split between authentic technical optimization ($42.6\%$) and promotional eWOM inflation ($57.4\%$).

%% =========================================================================
%% 6. MANAGERIAL DECISION SUPPORT AND IMPLEMENTATION GUIDELINES
%% =========================================================================
\section{Managerial Decision Support and Industrial Implementation Guidelines}
\label{sec:managerial_dss}

We translate the empirical architecture into four prescriptive operational rules for Chief Technology Officers (CTOs) and Chief Commercial Officers (CCOs):
\begin{enumerate}
    \item \textbf{Decision Rule 1 (Cross-Layer SLA Decoupling):} Telecommunications executive dashboards must decouple customer service level agreements (SLAs) into two independent channels: \textit{Core Infrastructure SLA} (RF latency, packet loss, OCS charging) and \textit{Digital Gateway SLA} (app crash-free sessions, OTP latency). Penalizing network infrastructure divisions for sentiment collapses caused by frontend mobile app build crashes destroys cross-functional alignment.
    \item \textbf{Decision Rule 2 (Dual-Track Surge-Adjusted Triage):} When monitoring customer sentiment during commercial campaigns, operational dashboards must report Track A (Organic Functional NSS) alongside aggregate VoC metrics. Executive bonuses and engineering evaluations should be tethered strictly to Track A to prevent marketing promotions from masking latent technical debt.
    \item \textbf{Decision Rule 3 (Canary Release Gating):} Telecommunications application release pipelines must integrate automated app store sentiment scraping with canary deployments. If an initial $5\%$ phased rollout triggers a statistically significant negative login sentiment spike ($p < 0.01$), the automated deployment pipeline must freeze rollout and trigger hotfix rollbacks before broad market contagion occurs.
    \item \textbf{Decision Rule 4 (Regulatory In-App Prompt Transparency):} Telecommunications regulatory authorities should mandate that operators disclose whether app store feedback is organically gathered or generated via incentivized in-app rating dialogs, standardizing VoC telemetry across competing licensees.
\end{enumerate}


\subsection{Executive Decision Framework: Mapping Empirical Telemetry to Business Strategy}
\label{sec:executive_decision_framework}

Bridging technical diagnostic performance with enterprise executive decision-making requires mapping statistical telemetry directly to telecommunications commercial key performance indicators (KPIs):

\begin{enumerate}
    \item \textbf{Customer Churn Mitigation and Lifetime Value (CLV) Protection:} In utility telecommunications, subscriber acquisition cost (SAC) is estimated to be 3.5$\times$ to 5.0$\times$ greater than subscriber retention cost \citep{reichheld2003one,parasuraman2005esqual}. Our empirical findings demonstrate an odds ratio of $\text{OR} = 4.23$ ($p = 0.0003$), proving that software application crashes are over four times more likely to induce catastrophic sentiment collapse than core network degradation. Crucially, digital self-care app users exhibit approximately 28\% higher average revenue per user (ARPU) and engage in significantly more digital data-pack purchases than non-app subscribers. Consequently, unmitigated mobile application defects disproportionately threaten the highest-value subscriber tiers. By triggering automated early-warning circuit breakers during canary releases, the expert system shields high-CLV cohorts from fatal product regressions.
    
    \item \textbf{CapEx versus OpEx Allocation Optimization:} When aggregate VoC or Net Promoter Scores plunge, traditional telecom management often initiates high-cost cross-functional escalations: deploying physical drive-test teams, realigning base station antennas, and contracting external radio frequency (RF) consultants. These represent capital-intensive (CapEx) and operationally cumbersome interventions. By systematically decoupling digital gateway software crashes from physical radio access network (RAN) outages, the expert system routes incident tickets directly to frontend mobile engineering teams (OpEx), preventing misdirected field investments and reducing diagnostic resolution times from weeks to minutes.
\end{enumerate}

To translate these strategic priorities into automated enterprise operations, Table~\ref{tab:executive_matrix} formalizes the multi-tiered \textbf{Executive Operational Decision Matrix}.

\begin{table*}[t]
\centering
\caption{Executive Operational Decision Matrix: Mapping Real-Time VoC Telemetry to Industrial Triage Workflows}
\label{tab:executive_matrix}
\footnotesize
\setlength{\tabcolsep}{6pt}
\begin{tabularx}{\textwidth}{@{}lllllX@{}}
\toprule
Alert Level & Statistical Trigger & Subsystem Attribution & Primary Root Cause & Automated Operational Workflow & Business Impact / KPI \\
\midrule
\textbf{Level 0 (Normal)} & $p \ge 0.05, \Delta\text{NSS} > -2.0$ & None (Steady-State) & Baseline organic variance & Routine synthetic telemetry ingestion; no engineering escalation & Steady subscriber sentiment; maintenance of baseline SLA \\
\addlinespace
\textbf{Level 1 (Advisory)} & $0.01 \le p < 0.05, \tau > 3.0$ & App / Commercial & Minor build friction or UI change & Flag ticket in Agile backlog; monitor next 200 reviews without release freeze & Proactive defect mitigation before churn contagion \\
\addlinespace
\textbf{Level 2 (Velocity Alert)} & $\tau > 10.0, \Delta\text{NSS}_{\text{promo}} > +15$ & Marketing / Promo & Promotional surge / gamification & Divert commercial reviews to Track B; audit app store gating prompts & Decouple marketing noise from core engineering evaluation \\
\addlinespace
\textbf{Level 3 (Critical Circuit)} & $p < 0.01, \Delta\text{NSS} \le -10.0$ & Client App Software & Fatal build crash / auth failure & Immediate CI/CD canary freeze; automated rollback to previous stable APK; zero RF field dispatch & Prevention of severe churn in high-ARPU digital subscriber tier \\
\addlinespace
\textbf{Level 3 (Network RAN)} & $p < 0.01, \text{Net} \gg \text{App}$ & Core Physical Network & Base station outage / fiber cut & Escalate to Network Operations Center (NOC); initiate BTS azimuth and power diagnostic & Protection of voice/data continuity; prevention of regulatory fines \\
\bottomrule
\end{tabularx}
\end{table*}

%% =========================================================================
%% 7. DISCUSSION AND LIMITATIONS
%% =========================================================================
%% 7. DISCUSSION, THREATS TO VALIDITY, AND LIMITATIONS
%% =========================================================================
\section{Discussion, Threats to Validity, and Limitations}
\label{sec:limitations}

\subsection{Theoretical and Practical Contributions}
This research introduces a systematic diagnostic framework in telecommunications customer intelligence. While existing industry practices rely on lagging, expensive retrospective surveys (e.g., periodic CSAT or quarterly NPS phone interviews), the proposed expert system operationalizes continuous, real-time perceptual telemetry. By integrating sub-word code-mixed NLP normalization with statistical quality control, the system successfully bridges raw linguistic feedback with automated infrastructure triage.

\subsection{Threats to Validity}
In alignment with standard guidelines for empirical systems research \citep{wohlin2012experimentation}, we systematically assess four primary dimensions of validity:

\begin{enumerate}
    \item \textbf{Construct Validity:} Construct validity concerns whether the quantitative telemetry accurately reflects underlying consumer dissatisfaction. We addressed this threat by: (a) formulating the Net Sentiment Score (NSS) to reflect net customer polarity; (b) calibrating apparent prevalence via the classical Rogan--Gladen estimator to correct for imperfect classifier sensitivity ($\mathrm{Se} = 0.711$) and specificity ($\mathrm{Sp} = 0.942$); and (c) validating multi-class topic boundaries against independent double-blind human consensus ratings ($94.12\%$ agreement).
    
    \item \textbf{Internal Validity:} Internal validity relates to whether confounding factors or experimental biases could have distorted the observed results. Potential annotator subjectivity was controlled through a formal multi-annotator guideline, double-blind evaluation, and resolution of the Kappa Prevalence Paradox ($\kappa = 0.5762$ vs.\ $\text{PABAK} = 0.8650$). Furthermore, to eliminate data contamination, all text normalization, vocabulary extraction, and TF-IDF bounds were fitted strictly on training folds during cross-validation.
    
    \item \textbf{External Validity:} External validity concerns the generalizability of our findings beyond the focal dataset. While our corpus spans 655,649 longitudinal records over a decade across all three major telecom operators in Bangladesh, harvesting was conducted via the Android Google Play platform. While Android represents over 95\% of smartphones in Bangladesh \citep{btrc2026quarterly}, this introduces a \textit{digital divide} selection effect: rural, low-income, or elderly subscribers utilizing 2G feature phones for voice calls are excluded from app store telemetry. Thus, the system is generalizable specifically to the rapidly expanding digital and self-care smartphone subscriber demographic.
    
    \item \textbf{Conclusion Validity:} Conclusion validity addresses the statistical credibility of the derived relationships. The statistical association between App Login defects and acute sentiment shocks was evaluated using both uncorrected Pearson $\chi^2(1) = 12.83$ ($p = 0.0003$, $\text{OR} = 4.23$) and Yates' continuity-corrected $\chi^2(1) = 11.55$ ($p = 0.0007$, $\text{RR} = 1.74$). We explicitly emphasize that these statistics reflect associations within self-care app store telemetry; they do not imply physical causal dominance over nationwide telecommunications infrastructure.
\end{enumerate}

\subsection{Study Limitations}
Despite its practical efficacy, two operational limitations warrant acknowledgment:
\begin{enumerate}
    \item \textbf{Absence of Direct Internal OSS/BSS Telemetry Pairing:} Due to strict proprietary non-disclosure restrictions of commercial telecommunications carriers, ground-truth network telemetry (e.g., physical base transceiver station [BTS] alarm logs, call detail records [CDRs], drive-test signal-to-interference ratios) could not be directly paired with individual reviews. Instead, network status was inferred from perceptual customer telemetry and public regulatory event timelines.
    \item \textbf{Static Gold Benchmark vs.\ Non-Stationary Slang Drift:} While the sub-word character n-gram ensemble demonstrated remarkable resilience across a ten-year longitudinal horizon, internet slang and transliterated code-mixing undergo continuous socio-linguistic mutation. A static supervised model inevitably experiences gradual semantic degradation over time unless reinforced through active continuous learning.
\end{enumerate}

%% =========================================================================
%% 8. CONCLUSION AND FUTURE WORK
%% =========================================================================
\section{Conclusion and Future Work}
\label{sec:conclusion}

This paper designed, implemented, and empirically validated \textbf{VoC-Intelligence}, an intelligent decision support system for telecommunications customer intelligence, cross-layer failure diagnostics, and automated software release gating. By coupling sub-word multi-lingual NLP normalization with a calibrated soft-voting ensemble, statistical quality control, and formal expert inference rules, the architecture resolves the long-standing industry dilemma of untangling client-side application crashes from backend physical network latency.

Evaluated over a comprehensive ten-year repository of 655,649 reviews across Bangladesh's three nationwide mobile network operators, the system achieved 87.50\% cross-validation accuracy and 94.12\% consensus agreement with human ground truth, while maintaining an ultra-low inference latency of 2.22 ms/review (450 reviews/sec on standard CPU hardware). Empirical findings indicate that within customer self-care feedback channels, mobile software build defects exhibit an odds ratio of $\text{OR} = 4.23$ ($p = 0.0003$; Yates-corrected $\chi^2(1) = 11.55, p = 0.0007$; relative risk $\text{RR} = 1.74$) of heading acute sentiment shocks compared to non-software categories. Furthermore, the formalized canary release gating protocol (Algorithm 1) establishes a practical industrial blueprint for automated release rollbacks in telecommunications CI/CD pipelines.

Building upon these findings, we delineate a concrete four-pillar roadmap for future research:
\begin{enumerate}
    \item \textbf{Multi-Modal Telecommunications Ingestion:} Expanding the sensory ingestion layer to incorporate unstructured customer service call-center audio recordings (via automatic speech recognition), Interactive Voice Response (IVR) transcripts, and USSD transaction failure logs.
    \item \textbf{Continual Active Learning and Dynamic Slang Adaptation:} Integrating self-supervised concept drift detection with human-in-the-loop active learning queues triggered when model predictive entropy exceeds critical thresholds ($\mathcal{H}(Y|X) > \tau_H$), allowing the expert system to autonomously ingest emerging Bengali slang.
    \item \textbf{Causal Impact and Econometric Validation:} Applying quasi-experimental econometric models (e.g., Difference-in-Differences and Synthetic Control Methods) to quantify the exact financial value and churn rate reduction directly attributable to automated canary hotfix rollbacks across multi-million subscriber bases.
    \item \textbf{Edge-Native Micro-Diagnostic Agents:} Quantizing lightweight inference models to run locally on mobile client devices, enabling immediate on-device failure detection and telemetry capture even during complete network transmission failure.
\end{enumerate}

%% =========================================================================
%% REFERENCES
%% =========================================================================
\bibliographystyle{elsarticle-harv}
\bibliography{references}

\end{document}
